% Converted from README.md
\section{"Fuzz" testiranje}\label{fuzz-testiranje}

Jedna od tehnika pronalaženja grešaka u programu koja se pokazala korisnom u praksi je \emph{fuzz testiranje} (eng. \emph{fuzzing}). Prilikom takvog testiranja, programu se šalju neočekivani i loši ulazni podaci, u cilju otkrivanja ulaza koji izazivaju greške. Standardni fuzzer alati mutiraju postojeće testove ili generišu testove na osnovu definisane gramatike ili skupa pravila. Fuzz testiranje vođeno pokrivenošću kôda (eng. \emph{coverage guided fuzzing}) prati putanje izvršavanja i daje prednost ulazima koji povećavaju pokrivenost. Tako rade postojeći alati za fuzz testiranje, kao što su American Fuzzy Lop (AFL), LLVM libFuzzer i HonggFuzz.

U najjednostavnijim slučajevima potrebno je da izolujemo funkcionalnost koju bismo da testiramo, napišemo nekoliko pratećih linija kôda, kompajliramo primer i pokrenemo fuzzer. Fuzzer će izvršiti hiljade ili desetine hiljada testova po sekundi i sačuvaće one interesantne ulaze koji povećavaju pokrivenost ili izazivaju određeno ponašanje.

Prvi korak u upotrebi \emph{libFuzzer}-a nad bibliotekom je kreiranje funkcije koja se u literaturi zove \emph{fuzz target}. Ta funkcija prihvata niz bajtova i na njih primenjuje funkcije iz biblioteke koju testiramo.

\begin{minted}{cpp}
extern "C" int LLVMFuzzerTestOneInput(const uint8_t *Data, size_t Size) {
    DoSomethingInterestingWithMyAPI(Data, Size);
    return 0;  // vrednosti osim 0 i -1 su rezervisane
}
\end{minted}

Ova funkcija ne zavisi ni u kojoj meri od \emph{libFuzzer}-a, pa ju je moguće i poželjno koristiti i sa drugim alatima za fuzz testiranje npr. AFL i/ili Radamsa.

Važne činjenice o fuzz funkciji (fuzz target):

\begin{itemize}
\item
  Alat za fuzz testiranje će izvršiti ovu funkciju mnogo puta sa različitim ulazima u okviru istog procesa.
\item
  Mora tolerisati sve vrste ulaznih podataka (prazne, velike, loše formatirane, itd.)
\item
  Ni za jedan ulaz se ne sme pozivati \texttt{exit()}.
\item
  Može koristiti niti, ali idealno bi bilo da sve niti budu spojene na kraju funkcije.
\item
  Mora biti deterministička, koliko god je moguće. Nedeterminizam, kao što su nasumične odluke nezavisne od ulaznih bajtova, učiniće fuzz testiranje neefikasnim.
\item
  Mora biti brza. Pokušajte zaobići kubne i veće složenosti, logovanje ili preveliko trošenje memorije.
\item
  Idealno bi bilo da ne modifikuje nijedno globalno stanje (iako to nije strogo neophodno).
\item
  Uža fuzz funkcija je obično bolja. Ako funkcija može da parsira više formata ulaza, treba je podeliti na nekoliko fuzz funkcija, po jednu za svaki format.
\end{itemize}

\emph{LibFuzzer} koristi sanitizer za pokrivenost kôda kojim se instrumentuje kôd da bi fuzzer dobijao informacije o pokrivenosti koje usmeravaju dalje fuzz testiranje. Podešava se opcijama iz grupe:

\begin{minted}{text}
-fsanitize-coverage=...
\end{minted}

Tokom fuzz testiranja treba uključiti i sanitizer koji prijavljuje greške tokom izvršavanja:

\begin{itemize}
\item
  \textbf{AddressSanitizer (ASAN)} detektuje greške sa memorijom. Uključuje se opcijom:
\begin{minted}{text}
-fsanitize=address
\end{minted}
\item
  \textbf{UndefinedBehaviorSanitizer (UBSAN)} detektuje upotrebu C/C++ konstrukcija koje uzrokuju nedefinisano ponašanje. Tipično se pokreće tako da se program zaustavi posle prijave:
\begin{minted}{text}
-fsanitize=undefined -fno-sanitize-recover=undefined
\end{minted}
  Može se uključiti i samo pojedinačna UBSAN provera, na primer:
\begin{minted}{text}
-fsanitize=signed-integer-overflow -fno-sanitize-recover=undefined
\end{minted}
  UBSAN se može kombinovati sa ASAN-om u istom prevođenju.
\item
  \textbf{MemorySanitizer (MSAN)} detektuje upotrebu neinicijalizovanih vrednosti, odnosno kôd čije ponašanje zavisi od sadržaja neinicijalizovane memorijske lokacije. Uključuje se opcijom:
\begin{minted}{text}
-fsanitize=memory
\end{minted}
  MSAN se ne sme kombinovati sa drugim sanitizerima i treba ga koristiti u zasebnom prevođenju.
\end{itemize}

Za svaku prijavu treba sačuvati ulaz koji je izazvao grešku, komandu kojom je fuzzer pokrenut i prvi relevantan okvir steka iz sanitizer izveštaja. Time fuzz nalaz postaje običan regresioni test.

\subsection{Korpus}\label{korpus}

Fuzzeri, kao \emph{libFuzzer}, koji koriste pokrivenost kôda za generisanje testova, oslanjaju se na korpus prethodno datih ulaza. Korpus treba da bude odmerena kolekcija validnih i nevalidnih ulaznih podataka. Fuzzer nasumičnim mutacijama generiše nove testove iz trenutnog korpusa. Ako mutacija aktivira izvršavanje prethodno nepokrivene putanje u kôdu koji se testira, tada se mutacija čuva kao nov test u korpusu i služi za buduće varijacije.

\emph{LibFuzzer} može da radi bez inicijalnog test primera, ali će biti manje efikasan ako biblioteka koja se testira prihvata kompleksne i strukturirane ulaze. Korpus može da posluži i za proveru da li svi ranije pronađeni testovi i dalje rade bez problema kada se propuste kroz kôd.

Ako imamo veliki korpus, bilo da je dobijen fuzz testiranjem ili drugačije, možemo ga smanjiti upotrebom opcije \texttt{-merge=1} tako da zadržimo pokrivenost.

\begin{minted}{bash}
$ mkdir NEW_CORPUS_DIR  # Store minimized corpus here.
$ ./my_fuzzer -merge=1 NEW_CORPUS_DIR FULL_CORPUS_DIR
\end{minted}

Ista opcija se može koristiti da se doda još testova u postojeći korpus. Samo oni testovi koji povećavaju pokrivenost će biti dodati u prvi navedeni korpus.

\begin{minted}{bash}
$ ./my_fuzzer -merge=1 CURRENT_CORPUS_DIR NEW_POTENTIALLY_INTERESTING_INPUTS_DIR
\end{minted}

Da bi se pokrenuo fuzzer sa korpusom testova, prvo treba kreirati direktorijum sa inicijalnim test primerima i potom ga pokrenuti:

\begin{minted}{bash}
$ mkdir CORPUS_DIR
$ cp /some/input/samples/* CORPUS_DIR
$ ./my_fuzzer CORPUS_DIR  # -max_len=1000 -jobs=20 ...
\end{minted}

Novi testovi biće dodavani u direktorijum korpusa.

Podrazumevano ponašanje je da fuzzer nastavlja da radi beskonačno ili dok ne pronađe grešku. Svaki pad ili greška koju detektuje sanitizer biće prijavljeni, fuzzer zaustavljen, a test primer koji je prouzrokovao grešku biće sačuvan u datoteku (obično, \texttt{crash-\textless{}sha1\textgreater{}}, \texttt{leak-\textless{}sha1\textgreater{}} ili \texttt{timeout-\textless{}sha1\textgreater{}}).

\subimport{01_libFuzzer/}{index.tex}

\subimport{02_tmux_fuzz/}{index.tex}

\subimport{03_jazzer/}{index.tex}

\subimport{04_afl++/}{index.tex}
